\chapter{Terrain et objet}
\label{ch:terrain}
% BUDGET : 2 pages (cf. memoire/PLAN.md, section 4). Etait a 4 : le contexte n'est pas
% la contribution, et le budget a ete rendu aux chapitres 3 et 4.
% THESE DU CHAPITRE : sur ce territoire, l'allocation des terres est un probleme de choix
% sous contraintes serrees -- et l'outil cense le traiter n'etait pas exploitable en l'etat.

\minisommaire{%
  \ms{sec:territoire}{climat, sol, chlordécone, foncier : un problème d'allocation}%
  \ms{sec:mosaica}{ce que le modèle promet}%
  \ms{sec:gams}{trois raisons de le réécrire}}

\section{Pertinence d'un modèle d'allocation en Guadeloupe}
\label{sec:territoire}

Les possibilités d'allocation de cultures par parcelles sont restreintes par le climat, le sol,
la pollution au chlordécone, le foncier et le travail. Les aides orientent le reste.

Environ \SI{14200}{\ha} sont classés à risque de contamination par la chlordécone, dont près de
\SI{7300}{\ha} encore en usage agricole. C'est un phénomène irréversible à l'échelle humaine car
la molécule persiste dans les sols pour des siècles \parencite{daafchlordecone}. Cela interdit
les cultures dont l'organe récolté est en contact avec le sol.

C'est une règle d'éligibilité par parcelle.

On compte \num{\obsParcelles} parcelles pour \num{\obsExploitations} exploitations. Le
parcellaire est fragmenté donc on ne peut a priori pas raisonner en exploitation moyenne. Le
régime des \gls{gfa} impose contractuellement une part minimale de canne à sucre à une partie de
ce foncier.

C'est une contrainte de politique publique.

D'après le jeu de données, sur les \SI{\obsSurfaceCultivee}{\ha} cultivés, la canne occupe
\SI{\obsHaCS}{\ha}, la prairie \SI{\obsHaPN}{\ha} et la banane \SI{\obsHaBA}{\ha}. Ces trois
postes correspondent à \SI{90}{\percent} de la surface totale. Un développement du vivrier passe
nécessairement par leur diminution, au profit de la surface vivrière, pour l'instant résiduelle.

C'est une réallocation.

On peut arbitrer sous contrainte. Le programme d'optimisation est donc pertinent en tant que
modèle d'allocation en Guadeloupe.

\section{\glsentryshort{mosaica} : ce que le modèle promet}
\label{sec:mosaica}

\gls{mosaica} manipule une mosaïque, c'est-à-dire une affectation simultanée de chaque parcelle
à une culture, sur tout un territoire, sous les contraintes propres à chaque exploitation. On
peut créer des scénarios (ex : ajout d'une contrainte, baisse de prix, etc.) et les évaluer
selon des indicateurs calculables depuis la mosaïque finale (ex : la marge, l'emploi, l'azote
épandu, l'eau prélevée, etc.). On peut donc réaliser des évaluations \emph{ex ante}
(chapitre~\ref{ch:explorer}).

La formulation du modèle parle de parcelles, de cultures et d'exploitations mais jamais de la
Guadeloupe. En langage de développeur, on appelle ça le \enquote{soft coding}. C'est volontaire
et c'est un des principes importants du portage (§~\ref{sec:architecture}).

\textcite[§~2.6]{chopin2015} publie une procédure de comparaison entre l'assolement simulé et
l'assolement observé, avec ses seuils. Cela permet de scorer la qualité de la calibration du
modèle (chapitre~\ref{ch:calibrer}).

\section{Le \glsentryshort{gams} existant, et pourquoi il a fallu reconstruire}
\label{sec:gams}

Le modèle \gls{gams} pèse \num{646}~ko répartis en \num{12} fichiers. On distingue
\num{3} raisons qui justifient sa réécriture.

\begin{enumerate}[leftmargin=1.6em,itemsep=6pt]

\item \gls{gams} est un environnement propriétaire qui fonctionne par licences. Cela crée des
problèmes de compatibilité de versions et crée une dépendance au logiciel.

On préfère adopter une approche open-source. On choisit Python comme langage de programmation
pour plusieurs raisons.

\begin{itemize}[leftmargin=1.2em,itemsep=3pt]
  \item Beaucoup de personnes sont formés à Python. Cela peut évacuer la question de la
    formation à un nouveau langage à la fois pour des collaborations entre équipes ou bien lors
    de recrutements.
  \item La forte concurrence entre les bibliothèques Python et leur modèle open-source assurent
    un large choix de modules optimisés et maintenus grâce à des mises à jour gratuites
    régulières. On choisit la bibliothèque d'optimisation Pyomo \parencite{pyomo2021} qui propose
    le solveur non-commercial \gls{highs} \parencite{highs2018}.
  \item Python permet très efficacement de faire de l'analyse de données approfondie et fait
    facilement le lien avec des SIG (système d'information géographique).
  \item Python est parfaitement compatible avec l'utilisation de GitHub. Cela permet de
    simplifier la collaboration en proposant des outils de gestion de projet et en rendant le
    code accessible à des équipes extérieures.
\end{itemize}

\item Dans le modèle \gls{gams}, faire varier un levier suppose d'éditer le code lui-même.
Réaliser une campagne de scénarios suppose un nombre considérable de modifications manuelles. La
traçabilité est faible, bricolée au cas par cas.

\item % troisieme raison : a rediger a partir de la reponse ci-dessous.

\begin{questionreponse}{Explique mieux la troisième raison.}
La troisième raison n'est ni l'accès ni le confort : c'est que le \gls{gams} \textbf{ne signale
pas ses propres défauts}.

Interrogé sur une colonne absente d'une table, \gls{gams} ne renvoie ni erreur ni avertissement :
il renvoie zéro. L'équation \texttt{Eq\_VE\_PLUIE} teste ce zéro contre la valeur~1 pour décider
si une parcelle peut porter un verger pluvial. Le test $0 \neq 1$ étant toujours vrai, la
variante est interdite sur \emph{toutes} les parcelles du territoire.

Ce défaut est invisible des deux côtés à la fois : le code reste correct d'apparence, et le
résultat reste plausible --- un verger de moins dans un assolement guadeloupéen n'a l'air de
rien. \emph{Un paramètre absent ne laisse pas un trou, il laisse un nombre crédible.} Le
découvrir suppose donc de refaire le modèle contre ses données, ligne à ligne : c'est la seule
lecture qui confronte chaque paramètre à la table qui est censée le porter.

Trois défauts de cette famille ont été trouvés pendant le portage
(§~\ref{sec:parite}, annexe~\ref{ann:portage}). La branche \enquote{rester en \gls{gams} et
n'ajouter qu'une couche de scripts autour} a donc été fermée par une \emph{mesure} et non par une
préférence de langage : une couche de scripts n'aurait rien révélé, puisqu'elle laisse le
programme intact.
\end{questionreponse}

\end{enumerate}

\begin{reflexion}
REFLEXIVITE
\end{reflexion}
